\scanpage{178}

U N I V E R S I T E

DES SCIENCES ET TECHNIQUES

DU LANGUEDOC

\verb|^^^^^^^^^^^^|

\hfill Montpellier, le 14 octobre 1980

Service de la Scolarité

MA/gs

Diffusion : Liste A

N° 5255

\begin{center}
\uline{AVIS DE SOUTENANCE D'UNE THESE DE 3e CYCLE}

``SIGNAUX ET BRUITS - Option 2''

\uline{Traitement du Signal et Communication}
\end{center}

\uline{NOM et prénom du candidat} : AMEZIANE \unsure{Rajae|Rajab}

\uline{Titre de la thèse} : ``ACQUISITION OPTIMISEE DES DONNEES EXPERIMENTALES PROCEDURES ET PROGRAMMES''

\uline{Date, heure et lieu de soutenance prévu} : 17 octobre 1980 à 11 h.15

Salle des Séminaires - C.E.E.S. - Université des Sciences et Techniques du Languedoc

\uline{Nom du directeur de thèse} : M. FILLARD J.P.

\uline{Désignation du Laboratoire où la thèse est préparée} : Laboratoire de Physique des Solides III - U.S.T.L.

\uline{Analyse de l'ouvrage} : Après un rappel de la théorie des phénomènes thermo-\newline
stimulés et des utilisations des divers types de mesures (TSC, TSL, vidages frac-\newline
tionnés, etc...), on décrit les dispositifs cryogéniques utilisés dans les zones de températures adaptées (77K-400K et 4K-300K). On présente successive-\newline
ment les diverses méthodes et programmes automatiques de mesure et d'exploita-\newline
tion des données physiques ainsi que des programmes assurant le fonctionnement du matériel cryogénique dans un mode automatisé.

\uline{Durée approximative des recherches ayant abouti à la thèse} : 3 ans

\uline{Composition du jury proposé} :
\begin{tabular}{l}
M. SAVELLI M.\\
M. FILLARD J.P.\\
M. LECOY G.\\
M. CAMBON G.\\
M. CASTAGNE M.
\end{tabular}

\begin{flushright}
P/ le Secrétaire Général,\\
l'Attachée d'administration\\
universitaire,

\unsure{Gentil}

\uline{A. GENTIL}
\end{flushright}
\margin{178}
